\subsection{UNIFORMLY CONTINUOUS FUNCTIONS}

DEFINITION 1. A mapping \(f\) of a uniform space \(\mathbf{X}\) into a uniform space \(\mathbf{X}'\) is said to be uniformly continuous if, for each entourage \(\mathbf{V}'\) of \(\mathbf{X}'\) , there is an entourage \(\mathbf{V}\) of \(\mathbf{X}\) such that the relation \((x,y)\in \mathbf{V}\) implies \((f(x),f(y))\in \mathbf{V}'\) .

In more expressive terms we may say that \(f\) is uniformly continuous if \(f(x)\) and \(f(y)\) are as close to each other as we please whenever \(x\) and \(y\) are close enough.

If we put \(g = f \times f\) , then Definition 1 means that whenever \(V'\) is an entourage of \(X'\) , \(\overline{g}^1(V')\) is an entourage of \(X\) .

Examples. 1) The identity mapping of a uniform space onto itself is uniformly continuous.

\begin{enumerate}
\def\labelenumi{\arabic{enumi})}
\setcounter{enumi}{1}
\item
  A constant mapping of a uniform space into a uniform space is uniformly continuous.
\item
  Every mapping of a discrete uniform space into a uniform space is uniformly continuous.
\end{enumerate}

PROPOSITION 1. Every uniformly continuous mapping is continuous.

This is an immediate consequence of the definitions.

On the other hand, a continuous mapping of a uniform space X into a uniform space \(X'\) need not be uniformly continuous, * as is shown by the example \(x \to x^{3}\) , a homeomorphism of R onto itself, which is not uniformly continuous with respect to the additive uniformity. * (See § 4, no. 1, Theorem 2.)

PROPOSITION 2. (a) If \(f: \mathbf{X} \to \mathbf{X}'\) and \(g: \mathbf{X}' \to \mathbf{X}''\) are two uniformly continuous mappings, then \(g \circ f: \mathbf{X} \to \mathbf{X}''\) is uniformly continuous.

\begin{enumerate}
\def\labelenumi{(\alph{enumi})}
\setcounter{enumi}{1}
\tightlist
\item
  A bijection \(f\) of a uniform space \(\mathbf{X}\) onto a uniform space \(\mathbf{X}'\) is an isomorphism if and only if \(f\) and the inverse of \(f\) are uniformly continuous.
\end{enumerate}

This follows immediately from the interpretation of Definition 1 in terms of the product mapping \(f \times f\) .

